\section{Relational Database ও Key}

\begin{definitionbox}{}
সারি ও কলামের টেবিলে ডেটা সংরক্ষণ এবং কী-এর সাহায্যে সংশ্লিষ্ট ডেটা যুক্ত করার নীতির ওপর গঠিত ডেটাবেজকে \textbf{রিলেশনাল ডেটাবেজ} বলে।
\end{definitionbox}


\subsubsection{গ্রাহক ও ক্রয়ের তথ্যের Relational বিন্যাস}
একজন গ্রাহক যতবার পণ্য কেনেন, প্রতিবার তাঁর নাম, ফোন নম্বর ও পরিচয়ের সব তথ্য নতুন করে লিখলে
একই তথ্যের অনেক কপি তৈরি হয়। ফোন নম্বর পরিবর্তনের সময় কোনো একটি কপি সংশোধন করতে ভুলে গেলে
একই গ্রাহকের দুই অর্ডারে দুই নম্বর দেখা যেতে পারে।

Customer টেবিলে গ্রাহকের পরিচয় এবং Orders টেবিলে প্রতিটি ক্রয়ের তথ্য আলাদা রেখে CustomerID দিয়ে যুক্ত করা যায়।
এতে গ্রাহকের বর্তমান তথ্য একটি নির্ধারিত উৎসে থাকে; প্রয়োজন হলে কুয়েরির সাহায্যে সেই তথ্য অর্ডারের সঙ্গে
মিলিয়ে পাওয়া যায়। ব্যাংকে Customer--Account এবং কলেজে Student--Payment সংযোগেও একই নীতি ব্যবহার করা হয়।

তবে কোনো অর্ডারের সময় পণ্য পৌঁছে দেওয়ার ঠিকানা (delivery address) সংরক্ষণের প্রয়োজন হলে তা সেই অর্ডারের
নিজস্ব তথ্য হিসেবে রাখতে হবে। গ্রাহকের বর্তমান ঠিকানা বদলালে পুরোনো অর্ডারের delivery address বদলানো উচিত নয়।


\begin{definitionbox}{RDBMS}
রিলেশনাল মডেলের টেবিল, কী ও শুদ্ধতার বিধি অনুসারে ডেটাবেজ তৈরি ও পরিচালনার সফটওয়্যারকে Relational Database Management System (RDBMS) বলে।
\end{definitionbox}
MySQL, PostgreSQL, Oracle Database, Microsoft Access ও SQLite পরিচিত রিলেশনাল ব্যবস্থা।

\subsection{Relational মডেলের বৈশিষ্ট্য}
\begin{enumerate}
\item \textbf{কাঠামো:} টেবিল, সারি ও কলাম দিয়ে ডেটা সংগঠিত হয়। প্রতিটি কলাম একটি নির্দিষ্ট বৈশিষ্ট্য এবং প্রতিটি সারি একটি এন্টিটির তথ্য প্রকাশ করে।
\item \textbf{শুদ্ধতার বিধি:} বিধি দিয়ে বৈধ মান ও সম্পর্কের সামঞ্জস্য নির্ধারিত হয়। প্রতিটি ঘরে নির্ধারিত নিয়ম অনুযায়ী একটি মান থাকে।
\item \textbf{ডেটা পরিচালনা:} কুয়েরি বা নির্দেশ দিয়ে প্রয়োজনীয় তথ্য নির্বাচন ও পরিবর্তন করা হয়।
\item \textbf{রেকর্ডের পরিচয়:} সারির অবস্থান স্থায়ী পরিচয় নয়; কী দিয়ে রেকর্ড শনাক্ত করা হয়।
\end{enumerate}
\textbf{উদাহরণ:} শিক্ষার্থী টেবিলের কাঠামো নির্ধারণের পর StudentID অনন্য রাখা এবং GPA-এর সীমা নির্ধারণ করা যায়; পরে শর্ত দিয়ে নির্দিষ্ট শিক্ষার্থী নির্বাচন করা যায়।
RDBMS হলো DBMS-এর একটি বিশেষ শ্রেণি। সব RDBMS ডেটাবেজ পরিচালনা করে, তবে সব DBMS রিলেশনাল মডেল অনুসরণ করে না।


\begin{flushleft}\small
\begin{tabularx}{\linewidth}{|p{24mm}|Y|Y|}\hline\chSixHeader
বিবেচনা & DBMS & RDBMS\\\hline
পরিসর & ডেটাবেজ পরিচালনার সফটওয়্যারের সাধারণ শ্রেণি & রিলেশনাল মডেল অনুসরণকারী DBMS\\\hline
কাঠামো & ব্যবস্থাভেদে বিভিন্ন ডেটা মডেল & সারি--কলামের সারণি ও সম্পর্কিত বিধি\\\hline
সম্পর্ক & মডেল অনুযায়ী প্রকাশ পায় & কী ও সংশ্লিষ্ট কলাম দিয়ে সংযোগ\\\hline
নিরাপত্তা/বহু ব্যবহারকারী & ব্যবস্থাভেদে সমর্থিত হতে পারে & এখানেও সুবিধা ও সীমা পণ্যনির্ভর\\\hline
\end{tabularx}
\end{flushleft}
\textbf{নির্বাচনের যুক্তি:} ভর্তি, ফি ও পরীক্ষার ফলের নির্ধারিত কাঠামো এবং তাদের নির্ভরযোগ্য সংযোগে RDBMS উপযোগী। সঠিক কী ও নিয়ম তথ্য মেলাতে সাহায্য করে; কিন্তু ভুল সম্পর্ক, অদক্ষ কুয়েরি বা অপরিকল্পিত কাঠামো থাকলে কেবল RDBMS ব্যবহারেই সমস্যার সমাধান হয় না।

\subsubsection{Relation ও SQL Table-এর বৈশিষ্ট্য}
\begin{flushleft}\small
\begin{tabularx}{\linewidth}{|p{27mm}|Y|Y|}\hline\chSixHeader
বিবেচনা & গাণিতিক Relation & বাস্তব SQL Table / ফল\\\hline
সারির পুনরাবৃত্তি & স্বতন্ত্র সারির সমষ্টি; সম্পূর্ণ একই সারি দ্বিতীয়বার থাকে না & উপযুক্ত Key বা UNIQUE বিধি না থাকলে একই মানের সারি যোগ হতে পারে\\\hline
অনন্যতা & রিলেশনের গঠনের অংশ & প্রয়োজনীয় বিধি কার্যকর করতে হয়; SELECT-এর ফলেও পুনরাবৃত্ত সারি থাকতে পারে\\\hline
\end{tabularx}
\end{flushleft}
\textbf{মনে রাখবে:} রিলেশনাল মডেলের নীতি এবং সফটওয়্যারে সেই নীতি কার্যকর করার বিধি আলাদা বিষয়।


\begin{enumerate}
\item \textbf{কলামের পৃথক পরিচয়:} একই টেবিলে কলামের নাম পৃথক থাকে। Name ও GPA ভিন্ন Attribute।
\item \textbf{মানের নির্ধারিত ধরন:} GPA-তে GPA-এর উপযোগী মান থাকে; নামের জায়গায় ইচ্ছামতো নম্বর বসানো যায় না।
\item \textbf{স্থায়ী পরিচিতি:} ওপরের সারি হওয়া পরিচয় নয়; StudentID দিয়ে শিক্ষার্থী শনাক্ত করা হয়।
\item \textbf{অনন্যতার প্রয়োগ:} নির্বাচিত Primary Key-তে একই পরিচয় দ্বিতীয়বার গ্রহণ করা হয় না।
\end{enumerate}


\subsection{Data Model}
\begin{definitionbox}{}
ডেটা কীভাবে সংগঠিত হবে, কী সম্পর্ক থাকবে এবং কী নিয়মে ব্যবহার করা যাবে তার ধারণাগত কাঠামো হলো \textbf{ডেটা মডেল}।
\end{definitionbox}
বিভিন্ন মডেলে একই বাস্তব তথ্য ভিন্নভাবে সংগঠিত করা যায়।
\begin{flushleft}\small
\begin{tabularx}{\linewidth}{|p{28mm}|X|X|}\hline
\chSixHeader মডেল & মূল বিন্যাস & উদাহরণ ও বৈশিষ্ট্য\\\hline
হায়ারার্কিক্যাল & গাছের মতো স্তর; সন্তানের একটি প্যারেন্ট & প্রতিষ্ঠান থেকে শাখা ও উপশাখা; নির্দিষ্ট স্তরবিন্যাস\\\hline
নেটওয়ার্ক & সংযুক্ত রেকর্ডের জাল & একই সদস্যের একাধিক প্যারেন্ট বা সংযোগ থাকতে পারে\\\hline
রিলেশনাল & সারি-কলামের টেবিল & পরিচিতি কী দিয়ে সংশ্লিষ্ট টেবিল যুক্ত করা যায়\\\hline
অবজেক্ট ওরিয়েন্টেড & অবজেক্টভিত্তিক উপস্থাপন & ডেটার সঙ্গে আচরণ বা পদ্ধতির ধারণা থাকে\\\hline
\end{tabularx}
\end{flushleft}
কলেজে শিক্ষার্থী ও বিভাগ আলাদা টেবিলে রাখা যায়। Student টেবিলের DeptID দিয়ে সংশ্লিষ্ট বিভাগের রেকর্ড চেনা যায়। সম্পর্কের ধরন ও বাস্তবায়ন \ref{sec:ch06-relationships} অংশে ব্যাখ্যা করা হয়েছে।


\subsubsection{Data Model-এর চিত্রভিত্তিক উদাহরণ}
হায়ারার্কিক্যাল মডেলে একটি মূল রেকর্ড থেকে নিচের স্তরে শাখা তৈরি হয়। যেমন কলেজের অধীনে বিভাগ এবং বিভাগের অধীনে শিক্ষার্থী দেখানো যায়। এই মডেলের নির্দিষ্ট প্যারেন্ট-চাইল্ড কাঠামো সহজ, কিন্তু একই শিক্ষার্থীকে একাধিক স্বতন্ত্র বিভাগের সঙ্গে যুক্ত করার প্রয়োজন হলে কাঠামো জটিল হতে পারে। নেটওয়ার্ক মডেলে এক সদস্যের একাধিক প্যারেন্ট-সংযোগের সুযোগ থাকায় এমন তথ্য অন্যভাবে প্রকাশ করা যায়।
\begin{flushleft}\centering\captionsetup{justification=centering}\begin{minipage}{\linewidth}\raggedright\small
\begin{tikzpicture}[every node/.style={booknode,font=\small},node distance=7mm]
\node[text width=25mm] (c) {কলেজ};
\node[below left=9mm and 11mm of c,text width=27mm] (s) {বিজ্ঞান বিভাগ};
\node[below right=9mm and 11mm of c,text width=27mm] (h) {মানবিক বিভাগ};
\node[below=of s,text width=27mm] (s1) {শিক্ষার্থী A};
\node[below=of h,text width=27mm] (h1) {শিক্ষার্থী B};
\draw(c)--(s);\draw(c)--(h);\draw(s)--(s1);\draw(h)--(h1);
\end{tikzpicture}
\captionof{figure}{স্তরবিন্যাসের একটি সরল উদাহরণ}
\end{minipage}\end{flushleft}
রিলেশনাল মডেলে একই বিষয় College, Department ও Student-এর মতো টেবিল ও কী দিয়ে প্রকাশ করা যায়। কোন মডেল বেছে নেওয়া হবে তা সংরক্ষণ ও ব্যবহারের প্রয়োজনের সঙ্গে সম্পর্কিত; শুধু দেখতে চিত্রটি সহজ কি না তা একমাত্র বিবেচনা নয়।

\textbf{জেনে রাখো—অন্য বিন্যাস:}
\textbf{List Structure} এখানে সরল তালিকাবিন্যাস বোঝাচ্ছে; এটি উপরের প্রধান formal database model-গুলোর সমতুল্য শ্রেণি নয়। এই বিন্যাসে তথ্য একটি সরল তালিকার মতো সাজানো থাকে। যেমন ধারাবাহিকভাবে সদস্যের নম্বর ও নাম লেখা তালিকা থেকে একেকটি এন্ট্রি পড়া যায়। এর সঙ্গে হায়ারার্কিক্যাল মডেলের প্যারেন্ট-চাইল্ড স্তরবিন্যাস এক নয়। \textbf{Multidimensional Model}-এ একাধিক মাত্রা ধরে তথ্য বিশ্লেষণ করা হয়। যেমন পণ্য, শাখা ও মাস---এই তিন দিক থেকে বিক্রয় দেখলে একই তথ্যের বিভিন্ন সারাংশ পাওয়া যায়। এখানে ``মাত্রা'' বলতে বিশ্লেষণের দিক বোঝায়; টেবিলের সারি-কলামের সংখ্যা নয়।

\subsubsection{Key ব্যবহারের প্রয়োজনীয়তা}
\begin{enumerate}
\item \textbf{সঠিক রেকর্ড শনাক্ত:} একই নামের দুই গ্রাহকের মধ্যে শুধু Name দেখে ফোন নম্বর বদলালে ভুল ব্যক্তির তথ্য সংশোধিত হতে পারে। CustomerID দিয়ে একজনকে নির্দিষ্ট করা পরিচয়ের সমস্যা সমাধান করে।
\item \textbf{তথ্যের সংযোগ:} গ্রাহকের একাধিক অর্ডারে তাঁর ID থাকা সংযোগের প্রয়োজন মেটায়।
\end{enumerate}
\textbf{মনে রাখবে:} ID দিয়ে খোঁজা গেলেই ব্যবহারকারী তা পরিবর্তনের অধিকার পান না—কী ও অনুমতি আলাদা বিষয়।

\subsection{Key-এর ধারণা ও প্রকারভেদ}
রেকর্ড শনাক্ত বা টেবিলের সংযোগ নির্ধারণে ব্যবহৃত একটি ফিল্ড বা একাধিক ফিল্ডের সমন্বয়কে \textbf{Key} বলা হয়। একটি নাম অনেক ব্যক্তির হতে পারে, কিন্তু প্রতিষ্ঠানের নিয়মে StudentID প্রত্যেক শিক্ষার্থীর জন্য পৃথক। নিচের উদাহরণে StudentID ও RegistrationNo উভয়ই বাধ্যতামূলক এবং প্রতিষ্ঠানে অনন্য ধরা হয়েছে। সারণিতে \textbf{Primary Key}-এর সংক্ষিপ্ত রূপ \textbf{PK} এবং \textbf{Foreign Key}-এর সংক্ষিপ্ত রূপ \textbf{FK} ব্যবহার করা হয়েছে।
\par\noindent\begin{minipage}{\linewidth}
\tcbset{chSixBox/.append style={breakable=false}}
\begin{definitionbox}{Primary Key (প্রাইমারি কী)}
একটি টেবিলের প্রতিটি রেকর্ডকে অনন্যভাবে শনাক্ত করার জন্য নির্বাচিত ফিল্ড বা ফিল্ডসমষ্টিকে প্রাইমারি কী বলে।
\end{definitionbox}
\noindent\begin{minipage}[t]{0.72\linewidth}\vspace{0pt}\raggedright\small
\begin{tabular}{|l|l|l|l|}\hline
\chSixHeader StudentID (PK) & Name & RegistrationNo & DeptID (FK)\\\hline
101 & Rupa & R501 & D01\\\hline
102 & Nabil & R502 & D01\\\hline
103 & Rupa & R503 & D02\\\hline
\end{tabular}
\end{minipage}\hfill\begin{minipage}[t]{0.25\linewidth}\vspace{0pt}
\begin{boardbox}{মনে রাখবে}
Name একই হতে পারে। নির্বাচিত PK একই বা NULL হতে পারে না। FK-এর পুনরাবৃত্তি নিজে ভুল নয়।
\end{boardbox}\end{minipage}\par
\end{minipage}\par



 এটি রেকর্ডকে অনন্যভাবে শনাক্ত করে এবং এর কোনো অংশ NULL হতে পারে না। একটি টেবিলে একটি প্রাইমারি কী নির্ধারিত হয়, তবে সেই কী একাধিক ফিল্ড নিয়ে গঠিত হতে পারে। এখানে StudentID ও RegistrationNo উভয়ই রেকর্ড আলাদা করতে পারে; প্রধান পরিচিতি হিসেবে StudentID নির্বাচন করা হলো।

\begin{definitionbox}{Composite Key}
একাধিক ফিল্ডের সমন্বয়ে গঠিত কী-কে \textbf{Composite Key} বলে।
\end{definitionbox}
 একই শ্রেণির প্রত্যেক শিক্ষার্থীর Roll আলাদা হলেও XI ও XII উভয় শ্রেণিতে Roll 1 থাকতে পারে। এই নিয়মে ClassName ও Roll একত্রে শিক্ষার্থীকে শনাক্ত করে। একাধিক শিক্ষাবর্ষের তথ্য রাখলে Session-ও পরিচিতির অংশ হতে পারে। নমুনায় মানগুলো আলাদা দেখালেই কোনো ফিল্ড সবসময় অনন্য হবে না; কী নির্বাচনে প্রতিষ্ঠানের নিয়মও বিবেচনা করতে হয়।

\begin{definitionbox}{Foreign Key (ফরেন কী)}
অন্য টেবিলের বা একই টেবিলের সংশ্লিষ্ট রেকর্ডের অনন্য কী-কে নির্দেশকারী ফিল্ড বা ফিল্ডসমষ্টিকে ফরেন কী বলে।
\end{definitionbox}
 Student-এর DeptID যদি Department-এর অনন্য DeptID-কে নির্দেশ করে, তবে Student.DeptID ফরেন কী। এক বিভাগের একাধিক শিক্ষার্থী থাকায় এই মান Student টেবিলে পুনরাবৃত্ত হতে পারে। ফরেন কী NULL হবে কি না তা নির্ধারিত নিয়ম ও NOT NULL বিধির ওপর নির্ভর করে।
\begin{flushleft}\small
\begin{tabularx}{\linewidth}{|p{29mm}|X|X|}\hline
\chSixHeader বিবেচ্য বিষয় & প্রাইমারি কী & ফরেন কী\\\hline
মূল কাজ & নিজ টেবিলের রেকর্ড শনাক্ত & উল্লেখিত রেকর্ডের সঙ্গে সংযোগ\\\hline
পুনরাবৃত্তি & সম্পূর্ণ কী-এর মান অনন্য & সাধারণভাবে পুনরাবৃত্ত হতে পারে\\\hline
NULL & গ্রহণযোগ্য নয় & নকশা ও বিধি অনুযায়ী\\\hline
সংখ্যা & একটি প্রধান কী; এক/একাধিক ফিল্ড & এক টেবিলে একাধিক ফরেন কী থাকতে পারে\\\hline
\end{tabularx}
\end{flushleft}

\textbf{কী নির্বাচনের উদাহরণ:} Name-এ Rupa দুইবার থাকায় তা প্রাইমারি কী নয়। DeptID-ও অনন্য নয়। StudentID স্থিতিশীল ও অনন্য হওয়ায় উপযুক্ত। কোনো টেবিলে এমন একটি ফিল্ড না থাকলে প্রতিষ্ঠানের নিয়মমাফিক যৌগিক কী বা আলাদা পরিচিতি নম্বর তৈরি করতে হয়।

\begin{boardbox}{মনে রাখবে}
প্রাইমারি কী-এর সম্পূর্ণ মান অনন্য এবং কোনো অংশ NULL নয়। কম্পোজিট প্রাইমারি কী-তে আলাদা অংশ পুনরাবৃত্ত হতে পারে; সম্পূর্ণ সমন্বয় নয়। এক-থেকে-অনেক (1:M) সম্পর্কে উল্লেখিত টেবিলের (Parent Table) একই রেকর্ডকে নির্দেশ করতে
উল্লেখকারী টেবিলে (Child Table) ফরেন কী-এর একই মান বারবার থাকতে পারে।
\end{boardbox}
\subsection{Relational মডেলের প্রবর্তন}
Edgar F. Codd ১৯৭০ সালে Relational Model প্রস্তাব করেন। ১৯৮৫ সালে তাঁর \textit{Codd's Twelve Rules} রিলেশনাল ব্যবস্থার বৈশিষ্ট্য ব্যাখ্যায় পরিচিতি পায়। মডেলের কাঠামো, শুদ্ধতার বিধি ও ডেটা পরিচালনার দিকগুলো আগের অংশে ব্যাখ্যা করা হয়েছে।

\subsection{গুরুত্বপূর্ণ প্রশ্ন}
\textbf{অনুশীলন:} (১) একটি টেবিলের প্রাইমারি কীতে তিনটি ফিল্ড থাকতে পারে কি? (২) একই নামের দুই কর্মচারীকে কীভাবে পৃথক করবে? (৩) Enrolment(StudentID, CourseID)-এ প্রতি শিক্ষার্থী একই কোর্সে একবার ভর্তি হলে কী নির্বাচন করবে?\par



